\documentclass[10pt,a4paper]{article}

% ----------------- Paquetes -------------------%
\usepackage[T1]{fontenc}
\usepackage[utf8]{inputenc}
\usepackage[a4paper,margin=2.5cm]{geometry}
\usepackage{mathptmx}             
\usepackage{changepage} 
\usepackage{setspace}
\usepackage[spanish]{babel}
\usepackage[labelfont=bf,labelsep=space]{caption}
\usepackage{amssymb}
\usepackage{amsmath}
\usepackage{ebgaramond}
\usepackage{placeins}
\usepackage{etoolbox}
\usepackage{setspace}
\usepackage{parskip} 
\usepackage{tcolorbox}
\usepackage{needspace}
\usepackage{xparse}
\usepackage{ocgx2}
\usepackage{hyperref}
\usepackage{hypcap}


%%%%%%%%%%%%%%%%%%%%%%%%%%%%%%%%%%%%%%%%%%%%%%%%%%%%%%%%%%%%%%%%
% --- Fija primero tus valores globales "buenos" ---
\setstretch{1.2}                 % Interlineado global
\setlength{\parskip}{0.7\baselineskip}
\setlength{\parindent}{0pt}
\newlength{\GlobalParskip}
\setlength{\GlobalParskip}{\parskip}

\newcounter{eqnum}[subsection]
\renewcommand{\theeqnum}{\arabic{eqnum}}
\newcounter{alphaitem}
\newcounter{numitem}

%% ------------------- Recuadros----------------------%%
\newcommand{\puntosclave}[1]{%
  \begin{tcolorbox}[
    colframe=black,
    title={Puntos clave},
    before upper={\setlength{\parindent}{0.5cm}\setlength{\parskip}{0.5\baselineskip}}
  ]
  #1
  \end{tcolorbox}
}

\newcommand{\modificaciones}[1]{
  \begin{tcolorbox}[
    colback=white,
    colframe=red,
    title={Modificaciones a realizar},
    before upper={\setlength{\parindent}{0.5cm}\setlength{\parskip}{0.5\baselineskip}}
  ]
  #1
  \end{tcolorbox}
}

%% ------------------- Preguntas TEST----------------------%%
% Opciones: radio (single-choice)
\NewDocumentCommand{\OptRadio}{m m m}{%
  \noindent\CheckBox[radio,name=#1,value=#2,width=1.2ex,height=1.2ex]{}%
  \;\textbf{#2.}~#3\par
}

% Opciones: checkbox (multi-choice)
\NewDocumentCommand{\OptCheck}{m m}{%
  \noindent\CheckBox[name=#1,width=1.2ex,height=1.2ex,bordercolor={0 0 0}]{}%
  \;#2\par
}

% Pregunta single-choice (radio)
% Args: 1=id, 2=enunciado, 3=opciones, 4=solución+justificación, 5=imagen (o {}), 6=origen
\NewDocumentCommand{\PreguntaTestSingle}{m m m m m m}{%
  \par\medskip
  \noindent\textbf{#1.}~#2\par\smallskip

    % Imagen (si no es {})
    \IfBlankTF{#5}{}{%
      \begin{center}
        \includegraphics[width=0.75\linewidth]{#5}
      \end{center}
    }

    \begin{Form}
      #3
    \end{Form}

    \par\smallskip
    \switchocg{sol-#1}{\fbox{Mostrar / ocultar solución}}%
    \hfill{\footnotesize #6}\par\smallskip

    \begin{ocg}{Solución}{sol-#1}{0}
      \noindent #4
    \end{ocg}
  \par\medskip
}

% Pregunta multi-choice (checkboxes)
\NewDocumentCommand{\PreguntaTestMulti}{m m m m m m}{%
  \par\medskip
  \noindent\textbf{#1.}~#2\par\smallskip

    % Imagen (si no es {})
    \IfBlankTF{#5}{}{%
      \begin{center}
        \includegraphics[width=0.75\linewidth]{#5}
      \end{center}
    }
    \begin{Form}
      #3
    \end{Form}

    \par\smallskip
    \switchocg{sol-#1}{\fbox{Mostrar / ocultar solución}}%
    \hfill{\footnotesize #6}\par\smallskip

    \begin{ocg}{Solución}{sol-#1}{0}
      \noindent #4
    \end{ocg}
  \par\medskip
}

%% ------------------- Listas----------------------%%
    % --- Lista alfabética ---
    \newenvironment{la}
      {\begin{list}{\alph{alphaitem})}{%
          \usecounter{alphaitem}%
          \setlength{\leftmargin}{1cm}%
          \setlength{\parskip}{3pt}% 
          \setlength{\itemsep}{0pt}% ← espacio entre ítems
          \setlength{\parsep}{0pt}%  ← párrafos dentro de un ítem
      }}
      {\end{list}}
      
    % --- Lista numérica ---
    \newenvironment{lnum}
      {\begin{list}{\arabic{numitem}.}{%
          \usecounter{numitem}%
          \setlength{\leftmargin}{1cm}%
          \setlength{\parskip}{3pt}% 
          \setlength{\itemsep}{0pt}% ← espacio entre ítems
          \setlength{\parsep}{0pt}%  ← párrafos dentro de un ítem
      }}
      {\end{list}}
      
% ------------------- Imágenes----------------------%
    \graphicspath{{Img/}}% Rutas por defecto 
    % --- Imagen adaptativa ---
    % Registros de dimensión definidos una sola vez
    \newdimen\imgcap
    \newdimen\imgmin
    \newdimen\imgmax
    \newdimen\imgremain
    \newdimen\imgh
    \newdimen\imgneed
    
    \newcommand{\imagenfit}[3]{%
      \addvspace{10pt}% separación antes
      \par\begingroup
        % Parámetros de composición (asignaciones, no \newdimen)
        \imgcap=1\baselineskip            % alto estimado de título+fuente
        \imgmin=0.15\textheight
        \imgmax=0.15\textheight
        \imgremain=\pagegoal
        \ifdim\pagegoal=\maxdimen \imgremain=0pt\fi
        \advance\imgremain by -\pagetotal
        \advance\imgremain by -\imgcap
        % Decisión de altura destino
        \ifdim\imgremain<\imgmin
          \newpage
          \imgh=0.20\textheight
        \else
          \imgh=\imgremain
          \ifdim\imgh>\imgmax \imgh=\imgmax \fi
        \fi
        % Reservar espacio para evitar cortes del bloque completo
        \imgneed=\imgh \advance\imgneed by \imgcap
        \Needspace{\imgneed}
        % Cuerpo
        \noindent\begin{minipage}{\textwidth}
          \centering
          \captionof{figure}{\textemdash\ #2}\par\vspace{0.3em}% título arriba
          \includegraphics[width=\linewidth,height=\imgh,keepaspectratio]{\detokenize{#1}}\par
          \vspace{0.3em}\small\textit{Fuente: }#3% fuente abajo
        \end{minipage}%
      \endgroup
      \par\addvspace{10pt}%
    }

%------------ Bloque de RECUADRO ESPECIAL -------------%
    \newenvironment{specialblock}[1]{%
      \begin{quote} % crea sangría a izquierda y derecha
      \small % texto más pequeño
      \noindent\textbf{#1}\\[-0.3em] % título en negrita
      \rule{\linewidth}{0.4pt}\\[0.6em]}{
      \end{quote}}
      
%-------------------------- Portada -----------------------%
\newcommand{\oldtitle}[1]{\def\OldTitle{#1}}
\newcommand{\changerationale}[1]{\def\ChangeWhy{#1}}

\newcommand{\changebox}{%
\begin{tcolorbox}[title={Historial de cambios del tema}]
\textbf{Título anterior:} \OldTitle\par
\textbf{Motivo del cambio:} \ChangeWhy
\end{tcolorbox}
}

%-------------------------- Author Font -----------------------%
    \newcommand{\authorfont}[1]{{\fontfamily{EBGaramond-TLF}\selectfont #1}}

%-------------------------- Anexos -----------------------%
    \makeatletter
    \newcounter{anexo}
    \renewcommand{\theanexo}{\Roman{anexo}}
    \newcommand{\anexo}[1]{%
      \clearpage
      \phantomsection
      \refstepcounter{anexo}%
      \noindent\textbf{Anexo \theanexo.- }#1\par\medskip
      \addcontentsline{stc}{section}{Anexo \theanexo.- #1}%
    }
    \makeatother

    %-------------------------- Lista de figuras (personal) -----------------------%
\makeatletter
\newcommand{\MyListOfFigures}{%
  \clearpage
  \phantomsection
  {\centering\bfseries\Large Índice de figuras\par}%
  \medskip
  % Entrada en nuestro índice de contenido (stc)
  \addcontentsline{stc}{section}{Índice de figuras}%
  % Recuperación del listado (sin cabecera propia)
  \begingroup
    \parindent=0pt\parskip=2pt
    \@starttoc{lof}%
  \endgroup
}
\makeatother

%-------------------------- Bibliografía -----------------------%
\makeatletter
% Contador para las referencias
\newcounter{myref}
% Cabecera de la bibliografía, entra en el índice 'stc'
\newcommand{\MyBibliography}{%
  \clearpage
  \phantomsection
  {\centering\bfseries\Large Bibliografía y referencias\par}%
  \medskip
  \addcontentsline{stc}{section}{Bibliografía y referencias}%
  \setcounter{myref}{0}%
}
\newcommand{\MyBibItem}[4]{%
  \refstepcounter{myref}%
  \noindent[\themyref]\ #1.\ \emph{#2}. #3. #4\par
}
\makeatother

%--------- Establecer cuatro niveles de índice --------%
    \setcounter{tocdepth}{4}   
    \setcounter{secnumdepth}{4}% numera hasta \paragraph

%%%%%%%%%%%%%%%%%%%%%%%%%%%%%%%%

\makeatletter

    %--------- Interlineado Global --------%
    \edef\GlobalBaseLineStretch{\baselinestretch}
    
    %--------- Espaciado de seccinones --------%
    \renewcommand\section{\@startsection{section}
      {1}{\z@}
      {2.5ex \@plus 1ex \@minus .2ex} % espacio antes
      {1.5ex \@plus .2ex}             % espacio después
      {\normalfont\Large\bfseries}    % formato del título
    }
    \renewcommand\subsection{\@startsection{subsection}{2}{\z@}
      {3ex \@plus 0.8ex \@minus .2ex}
      {1.5ex \@plus .2ex}
      {\normalfont\large\bfseries}
    }
    \renewcommand\subsubsection{\@startsection{subsubsection}{3}{\z@}
      {3ex \@plus 0.5ex \@minus .2ex}
      {1ex \@plus .2ex}
      {\normalfont\normalsize\bfseries}
    }
    \renewcommand\paragraph{\@startsection{paragraph}{4}{\z@}%
    {-2ex \@plus -0.5ex \@minus -.2ex}% espacio antes
    {0.8ex \@plus .2ex}% espacio después
    {\normalfont\normalsize\itshape}}% ← cursiva en lugar de negrita

    %--------- Ecuaciones --------%   
    % Variables globales para almacenar títulos y notas
    \gdef\leq@current@section{}%
    \gdef\leq@current@subsection{}%
    \gdef\leq@last@printed@section{}%
    \gdef\leq@last@printed@subsection{}%
    
    % Comando para añadir nota (invisible en el cuerpo)
    \newcommand{\eqnote}[1]{%
      \addcontentsline{leq}{leqnote}{#1}%
    }
    
    % Comando para ecuaciones
    \newcommand{\eqblock}[2]{%
      % Verificar si necesitamos imprimir el título de sección
      \ifx\leq@current@section\leq@last@printed@section\else
        \addcontentsline{leq}{leqsection}{\leq@current@section}%
        \global\let\leq@last@printed@section\leq@current@section
        \gdef\leq@last@printed@subsection{}% Reset subsección al cambiar sección
      \fi
      % Verificar si necesitamos imprimir el título de subsección
      \ifx\leq@current@subsection\leq@last@printed@subsection\else
        \ifx\leq@current@subsection\@empty\else
          \addcontentsline{leq}{leqsubsection}{\leq@current@subsection}%
          \global\let\leq@last@printed@subsection\leq@current@subsection
        \fi
      \fi
      % Ecuación
      \refstepcounter{eqnum}%
      \begin{equation*}
        #1\tag{E.\theeqnum}
      \end{equation*}%
      \addcontentsline{leq}{leqt}{[E.\theeqnum] -- #2}%
      \addcontentsline{leq}{leqm}{#1}%
    }
    
    %%% ----- Índice de ecuaciones ----------%%%
    % Tamaño para []
    \newcommand{\leqIndexFont}{\normalsize}
    \newcommand{\leqTitleFont}{\tiny}
    \newcommand{\leqNoteFont}{\tiny}
    % Cabecera de sección 
    \newcommand*\l@leqsection[2]{%
      \addvspace{25pt}% Mayor separación antes de nueva sección
      \noindent\leqIndexFont
      \makebox[\linewidth]{\rule{.38\linewidth}{0.4pt}\ \ \textbf{  \MakeUppercase{#1}}\ \ \rule{.38\linewidth}{0.4pt}}%
      \par\vspace{-10pt}
    }
    % Cabecera de subsección 
    \newcommand*\l@leqsubsection[2]{%
      \par\addvspace{15pt}%
      \noindent\leqIndexFont #1
      \par\vspace{-7pt}
    }
    % Título de ecuación 
    \newcommand*\l@leqt[2]{%
    \par\addvspace{10pt}%
      \par\noindent\leqTitleFont #1\par
    }
    % Miniatura de ecuación
    \newcommand*\l@leqm[2]{%
      \vspace{0.3\baselineskip}%
      \par\scriptsize\noindent\hspace*{1.5em}\(#1\)\par%
    }
    % Nota de ecuación 
    \newcommand*\l@leqnote[2]{%
      \vspace{0.2\baselineskip}%
      \par\noindent\leqNoteFont\hspace*{1.5em}\textbf{Nota.--} #1\par\vspace{0.5\baselineskip}%
    }
    % Generación del índice
    \newcommand{\mylistofequations}{%
      \clearpage
      \phantomsection
      % Título visual de la lista (ajústalo a tu gusto):
      {\centering\bfseries\Large Índice de ecuaciones\par}
      % Entrada en nuestro índice 'stc':
      \addcontentsline{stc}{section}{Índice de ecuaciones}%
      % Llamada a tu lista real (usa el nombre exacto de tu macro):
      \@starttoc{leq}%
    }
    
    %--------- Captura de títulos de sección --------%
    \let\leq@original@section\section
    \let\leq@original@subsection\subsection
    % Flag para saber si estamos en una sección asterisco
    \newif\ifLeqStarSection
    \LeqStarSectionfalse
    
    % Redefinir \section CON CONTROL DE ASTERISCO
    \renewcommand{\section}{%
      \@ifstar{\leq@section@star}{\@dblarg\leq@section@nostar}%
    }
    \def\leq@section@star#1{%
      \global\LeqStarSectiontrue
      \gdef\leq@current@section{#1}%
      \gdef\leq@current@subsection{}%
      \leq@original@section*{#1}%
      \global\LeqStarSectionfalse
    }
    \def\leq@section@nostar[#1]#2{%
      \global\LeqStarSectionfalse
      \gdef\leq@current@section{#2}%
      \gdef\leq@current@subsection{}%
      \leq@original@section[#1]{#2}%
    }
    % Redefinir \subsection
    \renewcommand{\subsection}{%
      \@ifstar{\leq@subsection@star}{\@dblarg\leq@subsection@nostar}%
    }
    \def\leq@subsection@star#1{%
      \global\LeqStarSectiontrue
      \gdef\leq@current@subsection{#1}%
      \leq@original@subsection*{#1}%
      \global\LeqStarSectionfalse
    }
    \def\leq@subsection@nostar[#1]#2{%
      \global\LeqStarSectionfalse
      \gdef\leq@current@subsection{#2}%
      \leq@original@subsection[#1]{#2}%
    }

    % ----- Restauración de valores Globales de estilo ----------%
    % Flags
    \newif\ifInFootnote
    \InFootnotefalse
    \pretocmd{\@footnotetext}{\InFootnotetrue}{}{}
    \apptocmd{\@footnotetext}{\InFootnotefalse}{}{}
    % Profundidad de listas (soporta anidadas)
    \newcount\MyListDepth \MyListDepth=0
    \AtBeginEnvironment{la}{\advance\MyListDepth by 1}
    \AfterEndEnvironment{la}{\advance\MyListDepth by -1}
    \AtBeginEnvironment{lnum}{\advance\MyListDepth by 1}
    \AfterEndEnvironment{lnum}{\advance\MyListDepth by -1}
    \AtBeginEnvironment{itemize}{\advance\MyListDepth by 1}
    \AfterEndEnvironment{itemize}{\advance\MyListDepth by -1}
    \AtBeginEnvironment{enumerate}{\advance\MyListDepth by 1}
    \AfterEndEnvironment{enumerate}{\advance\MyListDepth by -1}
    % Restauración diferida: hazlo al INICIO del próximo párrafo real
    \newif\if@pendingRestore
    \@pendingRestorefalse
    \newcommand{\RequestRestore}{\global\@pendingRestoretrue}
    \everypar{%
      \if@pendingRestore
        \global\@pendingRestorefalse
        \ifInFootnote\else
          \ifnum\MyListDepth=0
            \setlength{\parskip}{\GlobalParskip}%
            \setstretch{\GlobalBaseLineStretch}%
          \fi
        \fi
      \fi
    }
    % Al final de bloques:
    \AfterEndEnvironment{equation}{\RequestRestore}
    \AfterEndEnvironment{equation*}{\RequestRestore}
    \AfterEndEnvironment{la}{\RequestRestore}
    \AfterEndEnvironment{lnum}{\RequestRestore}
    % Al final de macros:
    \apptocmd{\eqblock}{\RequestRestore}{}{}
    \apptocmd{\imagenfit}{\RequestRestore}{}{}
    
\makeatother

% ===================== ToC propio con 4 niveles (único bloque) =====================
\makeatletter

% --- Escritores: añadimos una línea al .stc en cada cabecera numerada ---
% Guardamos las versiones actuales para llamar al comportamiento normal
\let\MyTOC@orig@section\section
\let\MyTOC@orig@subsection\subsection
\let\MyTOC@orig@subsubsection\subsubsection
\let\MyTOC@orig@paragraph\paragraph

% SECTION
\renewcommand{\section}{%
  \@ifstar{\MyTOC@section@star}{\@dblarg\MyTOC@section@nostar}%
}
\def\MyTOC@section@star#1{%
  \MyTOC@orig@section*{#1}% sin entrada al .stc
}
\def\MyTOC@section@nostar[#1]#2{%
  \MyTOC@orig@section[{#1}]{#2}% comportamiento normal (escribe al .toc)
  % Escribimos también nuestra línea al .stc (texto con \numberline)
  \addcontentsline{stc}{section}{\protect\numberline{\thesection}#2}%
}

% SUBSECTION
\renewcommand{\subsection}{%
  \@ifstar{\MyTOC@subsection@star}{\@dblarg\MyTOC@subsection@nostar}%
}
\def\MyTOC@subsection@star#1{%
  \MyTOC@orig@subsection*{#1}%
}
\def\MyTOC@subsection@nostar[#1]#2{%
  \MyTOC@orig@subsection[{#1}]{#2}%
  \addcontentsline{stc}{subsection}{\protect\numberline{\thesubsection}#2}%
}

% SUBSUBSECTION
\renewcommand{\subsubsection}{%
  \@ifstar{\MyTOC@subsubsection@star}{\@dblarg\MyTOC@subsubsection@nostar}%
}
\def\MyTOC@subsubsection@star#1{%
  \MyTOC@orig@subsubsection*{#1}%
}
\def\MyTOC@subsubsection@nostar[#1]#2{%
  \MyTOC@orig@subsubsection[{#1}]{#2}%
  \addcontentsline{stc}{subsubsection}{\protect\numberline{\thesubsubsection}#2}%
}

% PARAGRAPH
\renewcommand{\paragraph}{%
  \@ifstar{\MyTOC@paragraph@star}{\@dblarg\MyTOC@paragraph@nostar}%
}
\def\MyTOC@paragraph@star#1{%
  \MyTOC@orig@paragraph*{#1}%
}
\def\MyTOC@paragraph@nostar[#1]#2{%
  \MyTOC@orig@paragraph[{#1}]{#2}%
  % Si no numeras los \paragraph, cambia \theparagraph por texto plano sin \numberline
  \addcontentsline{stc}{paragraph}{\protect\numberline{\theparagraph}#2}%
}

% --- Impresor del índice propio (lee el .stc) ---
\newcommand{\MyTOC}{%
  \par\bigskip
  {\centering\bfseries\Large Índice de contenido\par}%
  \medskip
  \begingroup
    \parindent=0pt\parskip=2pt
    \@starttoc{stc}%
  \endgroup
}

\makeatother
% ===================== fin ToC propio =====================


%%%%%%%%%%%%%%


% --- Datos del tema ---
\title{Tema 3.A.43 -- Teorías del crecimiento económico (I)\\
\large Acumulación de capital y progreso técnico exógeno. El modelo de SOLOW. El modelo de SOLOW aumentado con capital humano. El modelo de RAMSEY-CASS-KOOPMANS}
\author{Marcelo Ortega}
\date{Diciembre 2025}
\newcommand{\TemaCode}{3A44}

\oldtitle{Tema 3.A.42. Teorías del crecimiento económico (I). El modelo de HARROD-DOMAR. El modelo de SOLOW. El modelo de RAMSEY-CASS-KOOPMANS. Sus implicaciones y limitaciones}
\changerationale{Se incluye explícitamente el modelo de Solow ampliado según el enfoque de MANKIW, ROMER y WEIL (1992). Se elimina el Modelo de HARROD-DOMAR, muy primitivo, para que haya suficiente tiempo para dedicarlo a los modelos más relevantes hoy. Ello no obstante, se hace notar que el modelo de Harrod introduce conceptos que siguen siendo de interés para los economistas del crecimiento económico, como las tasas de crecimiento garantizado y la tasa de crecimiento natural.}

\begin{document}


\maketitle
\changebox

\newpage

% --- Página de resumen y modificaciones ---
\puntosclave{ %% Escribir puntos clave Apuntes CECO
     Como punto de partida se toma el modelo AK, que pennite la presencia de crecimiento endógeno y crecimiento sostenido a largo plazo. Se recomienda dedicar unos 5 minutos a su exposición detallando la modelización de la tecnología, explicando el enfoque de Arrow sin entrar en detalle analítico, el equilibrio, la dinámica de la transición y las implicaciones. En cuanto al modelo AK con gasto público (modelo de Barro) se podrían dedicar unos 3 minutos a la exposición de la tecnología y los niveles óptimos de imposición y provisión de servicios públicos de consumo. A continuación, se presenta el modelo Lucas de capital humano. A efectos de la exposición oral sería aconsejable dedicar unos 5 minutos hacer hincapié en la acumulación de capital humano y en los detenninantes del crecimiento en estado estacionario. Respecto a la dinámica de la transición sería suficiente con mencionar las implicaciones.

     En la segunda parte del tema se presentan modelos con cambio tecnológico endógeno. En primer lugar, los modelos de variedad de productos de Romer. Para la exposición oral se recomienda dedicar unos 7-8 minutos explicando el modelo más simple e incorporando como conclusiones las implicaciones del modelo multisectorial. En segundo lugar, se analizan los modelos schumpeterianos de crecimiento endógeno, basados en la destrucción creativa. En la exposición oral se pueden emplear unos 7-8 minutos para analizar el modelo de un solo sector e incorporar las conclusiones del modelo con recursos naturales agotables, para subrayar la importancia del cambio tecnológico como motor del crecimiento económico.
    }
\vspace{1cm}

\modificaciones{
    
    }

\newpage
\MyTOC
\newpage

\mylistofequations
\clearpage

\MyListOfFigures
\clearpage


\pagenumbering{arabic}
\setcounter{page}{1}


\section{Introducción}\label{intro}
El objeto de este tema es el análisis de los modelos de crecimiento endógeno, que tratan de superar las limitaciones del modelo clásico, y se centran en el avance del conocimiento técnico para explicar el crecimiento económico. Por ello, es preciso conocer el tema 42, ya que se hará referencias a los modelos y desarrollos explicados en él. 

El modelo neoclásico de crecimiento explica el crecimiento a largo plazo en función de la tasa de cambio tecnológico que considera como dada exógenamente. Hay una buena razón, sin embargo, para pensar que el cambio tecnológico depende de decisiones económicas porque viene de innovaciones industriales realizadas por empresas que buscan maximizar beneficios y depende la financiación a la ciencia, de la acumulación de capital humano y de otras actividades económicas. La tecnología es, por tanto, una variable endógena, que debe ser determinada por el sistema económico. Las teorías del crecimiento deberían tener en cuenta esta endogeneidad, en tanto la tasa de progreso tecnológico es lo que determina la tasa de crecimiento a largo plazo. La nueva teoría del crecimiento desarrollada desde mediados de los años 80, se ha centrado en construir modelos con crecimiento sostenido que sean capaces de explicar, además, la experiencia de crecimiento de las economías y las diferencias en la renta per cápita de los países. Este desarrollo se ha realizado fundamentalmente en el contexto de modelos de crecimiento óptimo, recurriendo a la interpretación descentralizada de sus soluciones. Los modelos de la nueva teoría delcrecimiento consideran como endógeno el avance del conocimiento técnico, bien sea por laexistencia de un sector tecnológico que produce esos avances, bien sea por el aprendizaje que se logra con la construcción y uso de bienes de capital.

\subsection{Contextualización}\label{context}


\subsubsection{Relevancia}\label{importance}



\subsubsection{Problemática}\label{problem}



\subsection{Estructura de la exposición}\label{structure}



\clearpage
\section{Primera generación}
Incorporar la tecnología en la teoría del crecimiento nos obliga a tratar con el fenómeno de los rendimientos crecientes a escala.\footnote{En particular, los agentes deben tener incentivos a mejorar la tecnología, pero debido a que la función de producción agregada exhibe rendimientos constantes a escala en capital y trabajo, el teorema de Euler nos dice que todo el output de la economía será destinado a pagar a ambos factores por su producto marginal, sin que quede nada para pagar los recursos utilizados en mejorar la tecnología.
} Por tanto, una teoría de tecnología endógena no puede basarse en la teoría usual del equilibrio competitivo que requiere que todos los factores sean pagados por sus respectivas productividades marginales. 

La solución de ARROW (1962) a este problema era suponer que el progreso tecnológico era una consecuencia involuntaria de producir nuevos bienes de capital, un fenómeno denominado \textit{learning-by-doing}, que fue asumido como algo externo a las empresas. Es decir, si el progreso tecnológico depende de la producción agregada de capital y las empresas son todas muy pequeñas, entonces puede suponerse que tomarán la tasa de progreso tecnológico como dada independientemente de su propia producción de bienes de capital. Así cada empresa maximizará el beneficio pagando a sus factores: capital y trabajo, con sus respectivas productividades marginales sin ofrecer ningún pago adicional por su contribución al progreso tecnológico.

\subsection{El Modelo AK}
El modelo AK, introducido por REBELO (1991) se caracteriza por suponer una tecnología de rendimientos constantes a escala (tecnología AK) en el stock de capital. 

El hecho de imponer una relación lineal con el capital es lo que va a posibilitar la presencia de crecimiento endógeno ya que este incumplimiento de una de las condiciones de INADA abre la posibilidad de generar sendas de capital y de output per cápita con crecimiento sostenido en el largo plazo.

A continuación resolvemos el problema del planificador, que dada la ausencia de externalidades y bienes públicos, la solución competitiva será igual a la del planificador cumpliéndose los dos teoremas del bienestar.

\subsubsection{Supuestos}
El Modelo de REBELO parte de los sguientes supuestos:
\eqblock{
\begin{aligned}
    \begin{cases}
        \left.\begin{aligned}
            &\text{(1)}\quad \text{Mcdos. en CP; $P_{\not\perp}$; $S=I$}\\[0.5em]
            &\text{(2)}\quad \text{Tecnología $A$ constante}\\[0.5em]
            &\text{(3)}\quad \text{Tasa $S$ y $n$ constante}\\
            &\text{(4)}\quad \text{Función de Acumulación de $K$ agregado $\sim$ Inversión neta: }\dot K_t=\overbrace{I}^{sY}-\delta K_t
        \end{aligned}\right\}\text{\scriptsize$=$ SOLOW-SWAN}\\[0.5em]
        \text{(6)}\quad \underbrace{Y_t=AK_t}_{\text{\scriptsize Lineal en K}}\to y_t=Ak_t
        \begin{cases}
            \text{RCE}\\[0.5em]
            \not\downarrow PMg_K\to \text{No RMgD}\\[0.5em]
            \text{No INADA: }
            \begin{cases}
                \lim_{K\to0}PMg_K=A\neq\infty\\
                \lim_{K\to\infty}PMg_K=A\neq 0
            \end{cases}\\[0.5em]
        \end{cases}\\[0.5em]
    \end{cases}
\end{aligned}
}{Supuestos de partida. Modelo AK}

\subsubsection{Desarrollo analítico}
El planificador resuelve el siguiente problema de optimización en el cual se maximiza la utilidad intertemporal de los hogares sujeta a la restricción de recursos de la economía para cada instante temporal:

\eqblock{
\begin{aligned}
    &\boxed{\begin{aligned}
        &\max_{\{c_t\}}\int_0^\infty e^{-(\rho-n)t}\left(\frac{c_t^{1-\theta}-1}{1-\theta}\right)\mathrm{d}t\\[1em]
        &\text{s.a.}\quad
        \begin{cases}
            \dot k_t=Ak_t-(\delta+n)k_t-c_t\quad \text{\scriptsize (Ecuación dinámica)}\\[0.5em]
            k_0, \quad \text{dado}\\[0.5em]
            \lim_{t\to\infty}k_te^{-(A-\delta-n)t}=0
        \end{cases}
    \end{aligned}}\\[1em]
    &\left.\begin{aligned}
        &\substack{\text{\scriptsize Resolviendo por}\\\text{\scriptsize el Hamiltoniano}}\quad \text{CPO's}
        \begin{cases}
            c_t+\dot k_t-(n+\delta)k_t=Ak_t\\[1em]
            \underbrace{\frac{\dot c_t}{c_t}=\frac{1}{\theta}(A-\delta-\rho)}_{\text{\scriptsize Condición de EULER}}
        \end{cases}\\[1em]
        &\underbrace{\Longrightarrow}_{\text{\scriptsize Dividimos Ecuación Dinámica por $k_t$}}\quad \overbrace{\underbrace{\frac{\dot k_t}{k_t}=A-\frac{c_t}{k_t}-(\delta + n)}_{\text{\scriptsize Tasa de crecimiento capital per cápita}}}^{\text{\scriptsize \textbf{Constante}}}
        \end{aligned}\right\}
        \quad\Longrightarrow\quad 
        \begin{aligned}
            \underbrace{\boxed{\frac{\dot y_t}{y_t}=\frac{\dot k_t}{k_t}=\frac{\dot c_t}{c_t}=\frac{1}{\theta}(A-\delta-\rho)}}_{\text{\scriptsize Tasa de crecimiento ($\gamma$) constante para todos los parámetros ($i$)}}\\[0.5em]
            \underbrace{\sim}_{\substack{\text{\scriptsize Podemos }\\\text{\scriptsize aproximar}}}\quad \boxed{\gamma_i=sA-(\delta+n)}\\[0.5em]
            \text{Y, a nivel agregado: } \gamma^*=\gamma_i^*º+n
        \end{aligned}
\end{aligned}
}{Problema del planificador. Modelo AK (RABELO, 1990)}

Estas condiciones son las mismas que obteníamos en el equilibrio competitivo. Por tanto, podemos decir que el equilibrio competitivo es óptimo y que podemos encontrar unos precios (unos tipos de interés y un precio sombra o variable de coestado) tales que la solución del planificador coincida con el equilibrio competitivo. En definitiva, se satisfacen los dos teoremas del bienestar.

\paragraph{Representación gráfica: Tasa de crecimiento}
Una de las implicaciones más importantes del planteamiento/modelo AK es que, a diferencia del modelo neoclásico de crecimiento, puede generar crecimiento en términos per cápita a largo plazo en el producto interior bruto como observamos en la mayoría de los países del mundo, mediante crecimiento por acumulación (en esta caso, del stock de capital). Esta tasa de crecimiento será precisamente la diferencia entre la curva que venga representando la tasa de variación debida a la acumulación (CA) y la tasa de variación debido al crecimiento demográfico (CD).

\textbf{Incluir gráfico TASA de crecimiento (TECNOLOGÍA AK)}

\subsubsection{Valoración en términos de Política Económica}
Por tanto, los beneficios en términos de Política Económica son evidentes: es posible, de hecho real, inducir endógenamente el crecimiento económico (el incrmento gradual de la renta agregada en términos per cápita) mediante la acumulación de capital. Además, y en consecuencia, el ritmo de crecimiento vendrá definido por factores visibles, en particular: la tasa de ahorro ($s$), la variable tecnológica ($A$), el ritmo/tasa de crecimiento de la población ($n$) y el términos de depreciación del capital ($\delta$).

\paragraph{Implicaciones: Diferencias con el enfoque Neoclásico}
Podemos destacar las siguentes diferencias entre este modelo y el efoque neoclásico de crecimiento económico:
\begin{lnum}
    \item \textit{Crecimiento endogeneizado}. Puede haber una tasa de crecimiento de la renta/producción (PIB) per cápita positiva a largo plazo sin necesidad de tener que suponer que alguna variable crece continua y exógenamente;
    \item \textit{Equilibrio dinámico} (estacionariedad). La economía carece de una transición hacia el estado estacionario, ya que siempre crece a una tasa constante, por lo que siempre está en Estado Estacionario (siempre está creciendo a una tasa constante);
    \item \item \textit{No dinámica de transición}. La Economía carece de dinámica de transición que conduzca al Estado Estacionairo (es independiente de $k$);
    \item \textit{No convergencia}. El modelo predice que no existe ningún tipo de relación entre la tasa de crecimiento y el nivel de renta inicial. Dicho de otro modo, no predice ningún tipo de convergencia, ni condicional ni absoluta;
    \item \textit{Estructuralidad de los parámetros}. Cambios en los parámetros (shocks o política económica) tienen efectos permanentes sobre la tasa de crecimiento. Si hay una guerra que destruye nuestro capital y por tanto nuestra producción, luego seguiremos creciendo a la misma tasa que al principio, por lo que no nos recuperamos; y,
    \item \textit{Eficiencia dinámica}. Nunca existirá ineficiencia dinámica, es decir, una situación en la que haya demasiado ahorro.\footnote{Recordemos que para que esto ocurriese, en la zona dinámicamente ineficiente el tipo de interés en el estado estacionario era inferior a la tasa de crecimiento agregada, siendo el tipo de interés el producto marginal menos la tasa de depreciación. Aquí el tipo de interés es igual a $r = A - \delta$. Como la tasa de crecimiento per cápita es siempre igual a $[A-\rho-\delta]/\sigma$ la tasa de crecimiento agregada es $[A-\rho-\delta]/\sigma + n$. Para que haya ineficiencia dinámica es necesario que $A - \delta - n < [A-\rho-\delta]/\sigma$, lo que no puede darse ya que para que la utilidad esté acotada y el problema tenga sentido económico se tiene que dar la desigualdad al revés.
    
    Por tanto, el modelo nos demuestra que si al factor que se está acumulando se le elimina dicha productividad marginal decreciente, su acumulación puede ser motor de crecimiento sostenido.
    } Esto se debe a que una dinámica ineficiente implicaría $r^*>\gamma^*$. Por tanto, como $r^*=A-\delta$ y $\gamma^*=sA-(\delta + n)+n=sA-\delta$; como para encontrarse en la zona dinámicamente ineficiente debe cumplirse que $A-\delta<\underbrace{sA}_{\text{\scriptsize $s\in(0,1)$}}-\delta$, es imposible
\end{lnum}

\paragraph{Limitaciones}
Ahora bien, presenta importantes limitaciones que hay que destacar:
\begin{lnum}
    \item \textit{Ausencia de convergencia}. El modelo AK no predice convergencia condicional, pero en la evidencia empírica sí que se observa dicha convergencia. Esta limitación no se observa en el modelo de JONES y MANUELLI, quienes se basan en una función de producción que podría ser catalogada como intermedio entre SOLOW y REBELO que logra compatibilizar el crecimiento sostenido de la renta per cápita del modelo AK, con la posibilidad de convergencia condicional del modelo de SOLOW y que además, permite demostrar que la condición para lograr crecimiento sostenido sin progreso técnico no es que la productividad marginal del capital no sea decreciente, sino que no se cumpla una de las condiciones de Inada; y,
    \item \textit{Función de Producción ad hoc}. El modelo no nos da una verdadera explicación de cómo podría generarse esa función de producción tan características. REBELO asume que gracias al capital humano, pero en la práctica esto supondría que el capital humano y el capital físico son sustitutos perfectos que pueden agregarse alegremente. Lucas y Uzawa se apoyarán en el capital humano también pero como veremos con diferencias.
\end{lnum}

Por tanto, necesitamos teorías más sofisticadas que justifiquen hacer uso de una función de producción en la que el factor acumulado no tiene una productividad marginal decreciente. La literatura ha señalado la existencia de dos principales modelos.

\subsubsection{Extensiones: Externalidades}
Una vez expuesto el Modelo base y presentado tanto el concepto de \textit{tecnología AK} (es decir, tipo de función de producción que mantiene una relación lineal o positiva con un factor exógeno) como la vetaja que ofrece a la hora de fundamentar las tasas de crecimiento contsante observadas empíricamente, pasamos a ver las principales formas de modelizar este tipo de crecimiento endogeneizado. 

\paragraph{Externalidades de Capital ($K$)}
Existen dos formulaciones alternativas para introducir la existencia de externalidades positivas en el crecimiento generadas por el capital. No obstante, ambas parten de la siguiente proposición:

\textbf{Definición.-} (\textit{Externalidades dle capital}) La inversión en capital genera nuevos conocimientos que se proyectan como externalidades en dos vertientes: (i) como externalidades internas, la inversión en capital genera aprendizaje por la práctica ya que al invertir en éste se adquiere experiencia; y, (ii) como externalidad externa, ya que este conocimiento generado se desborda al resto de la economía, contagiéndose y propagándose de una empresa a otra sin necesariamente haber sido la que haya realizado la inversión. 

Una primera aproximación de introducir las externalidades del capital es la aproximación de LUCAS (1988), considerando que se genera cuando se produzca un incremento del capital per cápita:

\eqblock{
\begin{aligned}
    &\text{Función de Producción Agregada:}\quad \underbrace{Y=AK^\alpha L^{1-\alpha}\left(\frac{K}{L}\right)^\eta}_{\text{\scriptsize Externalidad cuando $\uparrow k$, capital per capita}}\\[1em]
    &\underbrace{\Longrightarrow}_{\text{\scriptsize En términos \textit{pc}}}\quad \boxed{y=Ak^{\alpha+\eta}}\\[2em]
    \substack{\text{Por}\\\text{tanto,}}\quad
    &\begin{cases}
        \text{Si, $\alpha+\eta<1$}\qquad \exists\text{Extern., pero pequeñas}\sim \overbrace{\text{Modelo SOLOW}}^{\text{\scriptsize CA decreciente en $k$}}\\[0.5em]
        \text{Si, $\alpha+\eta=1$}\qquad \exists\text{Extern., tecno AK}\to \text{Modelo ROMER}\sim\text{Modelo AK}\\[0.5em]
        \text{Si, $\alpha+\eta>1$}\qquad \exists\text{Extern. muy fuertes}\to \underbrace{\text{EE inestable}}_{\text{\scriptsize CA creciente en $k$}}
        \begin{cases}
            \text{Si $k<k^*_{ee}$}\to \text{Senda implosiva: $\dot k_t\to0$}\\[0.5em]
            \text{Si $k^*_{ee}<k$}\to \text{Senda explosiva: $\dot k_t\to\infty$}
        \end{cases}
    \end{cases}
\end{aligned}
}{Externalidades del capital (I). Modelo de LUCAS (1988)}

Por tanto, vemos como el Modelo de LUCAS nos permite introducir los efectos derivados de la incorporación de tecnología AK siempre y cuando el tamaño de la externalidad sea tal que el impacto de ésta sobre la producción/renta sea igual al de los trabajadores, es decir, $\alpha+\eta=1\to\eta=1-\alpha$.

Otra forma de introducir la presencia de externalidades del capital es la propuesta por ROMER (1986). En este caso asumiremos que la externalidad se genera cuando aumenta el stock de capital físico agregado (no per cápita):

\eqblock{
\begin{aligned}
    &\text{Función de Producción Agregada:}\quad \underbrace{Y=AK^{\alpha+\eta} L^{1-\alpha}}_{\text{\scriptsize Externalidad cuando $\uparrow K$, stock de capital}}\\[1em]
    &\underbrace{\Longrightarrow}_{\text{\scriptsize En términos \textit{pc}}}\quad \boxed{y=Ak^{\alpha+\eta}L^{\eta}}\\[2em]
    \substack{\text{Por}\\\text{tanto,}}\quad
    &\begin{cases}
        \text{Si, $\alpha+\eta=1$}\Longrightarrow\qquad \gamma_k=f(L)
        \Rightarrow\begin{cases}
            \exists\text{Efecto escala}\\[0.5em]
            \gamma_k^{\text{cte.}}\iff\eta=0
        \end{cases}
        \quad \neq\text{Evidencia empírica}
        \\[0.5em]
        \text{Si, $\alpha+\eta<1$}\Longrightarrow\qquad \exists\text{EE donde }k^*(L)\Longrightarrow\;\;(\uparrow L\Rightarrow\;\uparrow y)\quad \neq\;\text{Evidencia empírica}
    \end{cases}
\end{aligned}
}{Externalidades del capital (II). Modelo de ROMER (1986)}

Por tanto, comom vemos, las externalidades introducidas a través del capital agregado de la economía inducen efectos escala que son contradecidos (o no son contrastados) por la evidencia empírica.\footnote{No obstante, algunos autores defienden la existencia de este fenómeno, considerando que la unidad geográfico que utilizamos para comparar y contrastar los resultados empíricamente no tiene que ser necesariamente la unidad política nacional, sino que podría analizarse regionalmente, lo que permitiría o posibilitaría la existencia de dichos fenómenos.
}

\paragraph{Función de SOBELOW}
Otra forma de modelizar el crecimiento de la renta o producción agregada de manera endógena, sorteando las lmiitaciones del modelo de tecnología AK, es a través de la combinación (o integración) de la Función de Producción de REBELO y la de SOLOW.

\eqblock{
\begin{aligned}
    &\text{Función de Producción Agregada: }\quad Y_t=\underbrace{AK_t}_{\text{REBELO}}+\overbrace{BK_t^\alpha L_t^{1-\alpha}}^{\text{SOLOW}}\\[1em]
    &\underbrace{\Longrightarrow}_{\text{\scriptsize En términos pc}} \boxed{y_t=Ak_t+Bk_t^\alpha}\quad \text{Donde,}\;
    \begin{cases}
        \text{(1) Rendimientos Constantes a Escala}\\[0.5em]
        \text{(2) Rendimientos No decrecientes de $K$: $\not\downarrow PMg_K$}\\[0.5em]
        \text{(3) No INADA}\\[0.5em]
    \end{cases}\\[1em]
    &\text{Introduciendo ésta en la } \underbrace{\text{ecuación dinámica de SOLOW}}_{\dot k=sf(k)-(n+\delta)}:\quad \boxed{\overbrace{\frac{\dot k}{k}}^{\substack{\text{\scriptsize Dividimos por}\\\text{\scriptsize capital pc ($k$)}}}=\underbrace{sA+sBk^{\alpha-1}}_{CA}-\overbrace{(n+\delta)}^{CD}}\\[2em]
    &\substack{\text{\scriptsize Por }\\\text{\scriptsize tanto,}}
    \begin{cases}
        \text{(1) Asíntotas:}\quad 
        \begin{cases}
            CD\text{ (horizontal)}\\
            CA\text{ (decreciente)}
            \begin{cases}
                k\to0\Longrightarrow\;CA\to\infty\\
                k\to\infty\Longrightarrow\;CA\to sA\\
            \end{cases}
        \end{cases}\\[1em]
        \text{(2) Casos:}\quad 
        \begin{cases}
            \text{\textbf{{Si}}}\;\;\;n+\delta<sA\Longrightarrow\;CA\cap CD=\varnothing\Rightarrow
            \begin{cases}
                \gamma^*_{cp}>0\\
                \gamma^*_{lp}\to sA-(n+\delta)
            \end{cases}\quad (\sim\text{\scriptsize Tecnología AK})\\[0.5em]
            \text{Si}\;\;\;n+\delta>sA\Longrightarrow\;CA\cap CD\neq\varnothing\quad (\sim\text{\scriptsize SOLOW})
        \end{cases}
    \end{cases}
\end{aligned}
}{Función de SOBELOW. Modelo de tecnología AK}

Por tanto, como vemos, la principal implicación de este modelo es que es posible generar crecimieto endógeno aún en presencia de rendimientos decrecientes del factor productivo capital, siempre que no se cumplan las condiciones de INADA (lo que resulta más sencillo y posible de asumir, solventando los problemas/reticencias que generaba la construcción \textit{ad hoc} de la función de producción en el modelo AK). Además, nos permite intuir que, aunque los rendimientos de dicha función son importantes, son verdaderamente irrelaventes para explicar el crecimiento sostenido de la renta/producción a largo plazo (que se ve principalmente influenciada por las dinámicas de acumulación y crecimiento demográfico).

\subsection{Crecimiento por acumulación de capital humano} 
Hemos visto que tanto en el modelo baseline con tecnología AK, como en sus variantes, puede generarse crecimiento de forma endógena. Sin embargo, éstos aglutinan dentro de los componentes del capital físico tanto el propio capital como el capital humano generado, produciéndose estos con la misma tecnología (sin distinguir entre ambos). Tratando de superar estas limitaciones surge el modelo de crecimiento inducido por la acumulación de capital humano. 

\subsubsection{Supuestos}
Consideramos los siguientes supuestos de partida:

\eqblock{
\begin{aligned}
    \text{\scriptsize Supuestos:}\;
    \begin{cases}
        \text{(1) Existen dos sectores:}\;
        \begin{cases}
            Y\quad (K,CH)\longrightarrow\text{\scriptsize Producto final: Consumido o transformado en $K$}\\
            H\quad (CH)\longrightarrow\text{\scriptsize Educativo, produce y acumula $H$}
        \end{cases}\\[0.5em]
        \text{(2)}\;\exists A_1, A_2:\;
        \begin{cases}
            A_1\to Y\\
            A_2\to H
        \end{cases}\\[0.5em]
        \text{(3)}\;H\equiv\text{Bien rival:}\;
        \begin{cases}
            H_Y=\mu H\\
            H_H=(1-\mu)H
        \end{cases}
    \end{cases}
\end{aligned}
}{Supuestos. Modelo de LUCAS-UZAWA (acumulación de capital)}

Así, en este modelo se permite la existencia de dos sectores: el que produce el bien final y el que produce conocimiento (acumulación de capital humano). Suponemos que el sector que produce conocimiento, como dice la evidencia empírica, es más intensivo en capital humano y lo llevamos al limite asumiendo que el único input acumulable en la producción de conocimiento es el capital humano.

\subsubsection{Desarrollo}
Una vez que tenemos claros los supuestos, podemos plantear el problema de optimización de un planificador benevolente de la siguiente forma:\footnote{
Sobre el HAMILTONIANO, resolvemos las CPO's, tomamos logaritmos y derivamos respecto a $t$ para encontrar la senda óptima de consumo que se presenta.
}

\eqblock{
\begin{aligned}
    &\boxed{
    \begin{aligned}
        &\max_{\{c_t,\mu_t\}}U(0)=\int_{0}^\infty e^{-(\rho-n)t}\left(\frac{c_t^{1-\theta}}{1-\theta}\right)\mathrm{d}t\\[0.5em]
        &\text{s.a.}\;
        \begin{cases}
            \dot k_t=Ak^\alpha(\mu_t H)^{1-\alpha}-c_t-(\delta_K+n)k_t\\[0.5em]
            \dot h_t=B(1-\mu_t)h_t-(\delta_H+n)h_t\\[0.5em]
            k_0\;\;y\;\;h_0,\quad \text{dado}
        \end{cases}
    \end{aligned}}\quad \text{donde,}\;
    \begin{cases}
        B\equiv\quad \text{$\sim$ Productividad en el aprendizaje ($A_2$)}\\
        \text{Dos v. control:}\quad (c_t, \mu)\\
        \text{Dos v. estado:}\quad (k_t, h_t)\\
    \end{cases}\\[1em]
    &\text{(Hamiltoniano) CPO's}\Longrightarrow\text{Condición de Euler:}\underbrace{\quad\boxed{\frac{\dot c_t}{c_t}=\frac{1}{\theta}(\underbrace{\alpha k^{\alpha-1}_t(\mu_th_t)^{1-\alpha}}_{\text{Tecnología educación: $B$}}-(\delta+\rho))}}_{\text{Senda óptima de consumo}}\\[1em]
    &\Longrightarrow\quad \boxed{\gamma^*_{c}=\gamma^*_y=\gamma^*_k=\gamma^*_h=\frac{1}{\theta}(B-\delta-\rho)}
\end{aligned}
}{Desarrollo analítico. Modelo de LUCAS (capital humano)}

Por tanto, como vemos la tasa de crecimiento a largo plazo es similar a la obtenida en los modelos de tecnología AK. Sin embargo, el parámetro de productividad que afecta al crecimiento al largo plazo  de cualquier variable de la economía será, por el contrario, la productividad asociada al sector educativo ($B$), debido a que hemos asumido que el sector educativo no utiliza capital físico (y sí el sector de producción de bienes de consumo o bienes de capital). Por tanto, en el Estado Estacionario, la economía crece a una tasa constante mayor a cero siempre y cuando la productividad del sector educativo sea mayor a la tasa de depreciación de la economía y la impaciencia de los hogares en ésta. 

\subsubsection{Dinámicas de transición} 
A diferencia del modelo AK en el que la economía estaba siempre en una situación de estado estacionario, en este modelo hay un período de transición. 

Debido a que la tasa de crecimiento del consumo (y de la economía) se relaciona negativamente con la ratio $K/H$, el crecimiento tenderá a aumentar si el capital humano es relativamente abundante (en comparación con su nivel de estado estacionario) pero disminuirá si el capital humano es relativamente escaso.\footnote{Esto implica que ante una pérdida de capital físico (por ejemplo debido a una guerra), la economía se recuperará más rápidamente (capital físico escaso relativamente) que cuando hay una pérdida de capital humano (por una epidemia, por ejemplo), donde el crecimiento de la economía será más lento. Esta predicción se resume en que reponer capital físico es relativamente más fácil que reponer capital humano, hecho generalmente probado empíricamente.
} 

Así, hemos visto que en Estado Estacionario la ratio $k/h$ es constante, por lo que la dinámica de transición surge si $k/h$ es diferente de la ratio $k/h$ de estado estacionario. Es decir, la dinámica de transición surge por una descomposición entre los dos sectores:
\begin{la}
    \item \textit{Si $\frac{k_0}{h_0}<\frac{k^*_{ee}}{h^*_{ee}}$}. La tasa de crecimiento de la economía será superior durante el período de ajuste (Alemania, Japón, II Guera Mundial); y,
    \item \textit{Si $\frac{k_0}{h_0}>\frac{k^*_{ee}}{h^*_{ee}}$}. La tasa de crecimiento de la economía será inferior a la del estado estacionario durante el período de ajuste (pandemia). 
\end{la}

Este corolario es muy importante, pues recordemos que en el modelo AK, a diferencia de cualquiera de sus vaiaciones, cualquier shock dejaba a la economía desplazada para siempre pues no se recuperaba y mantenía el mismo crecimiento. 

\subsubsection{Valoración en términos de Política Económica}
La naturaleza del modelo de LUCAS hace que las conclusiones obtenidas sean similares a la de los modelos anteriores, al permitir que haya crecimiento sostenido al eliminar la productividad marginal decreciente del factor acumulado, en este caso el capital humano; así como por la ausencia de convergencia entre economías.

\paragraph{Implicaciones}
A corto largo plazo, el crecimiento del consumo, la acumulación de capital per cápita y el capital humano en términos per cápita dependerá de diversos parámetros:

\eqblock{
\begin{aligned}
    \boxed{\gamma^*_c=\gamma^*_y=\gamma^*_h=\gamma^*_k=f(\overbrace{B}^{\text{\scriptsize Positivamente}},\underbrace{\theta,\rho,\delta}_{\text{\scriptsize Negativamente}})}
\end{aligned}
}{Determinantes del crecimiento a largo plazo. Modelo de LUCAS}

La productividad del sector educativo es, en puridad, el verdadero motor de la economía, siempre y cuando la acumulación de éste no dependa también del capital físico. 

En términos de política económica esto tiene grande implicaciones, de las que cabe resaltar dos:
\begin{lnum}
    \item \textit{Importancia de la educación}. Debe prestarse especial atención a las políticas dirigidas a aumentar el capital humano (educación) de la sociedad puesto que éste constituye el motor del crecimiento futuro;
    \item \textit{Balance en ahorro y educación}. En el caso de que relajemos el supuesto límite de la función deproducción de capital humano, tanto la tecnología convencional como la productividad del sector educativo afectarán a la tasa de crecimiento a largo plazo y dependerá de la participación de cada uno en la producción.
\end{lnum}

En definitiva, este modelo supera las limitaciones subyacentes en la propuesta de REBELO de homogeneizar el capital de la función de producción. 

Dado que necesitamos que en el estado estacionario exista crecimiento no nulo, y que esto es posible si los rendimientos a escala no constantes en los inputs acumulables en ambos sectores, la tecnología que produce conocimientos será lineal en el capital humano. 

El análisis de este modelo teórico ha desarrollado una gran variedad de estudios empíricos relacionados con las preguntas:
\begin{lnum}
    \item \textit{¿Debería la educación ser financiada de forma privada o pública?}
    \item \textit{¿Pueden las políticas educativas incentivar el crecimiento económico?}
    \item De ser así, \textit{¿qué se debería enfatizar, la educación primaria/secundaria o la educación superior?}
    \item \textit{¿Debería el gobierno subsidiar la educación fomial o la formación profesional en el ámbito privado/empresarial?}
    \item \textit{¿Debería financiarse más la educación elitista o la generalista?}
    \item \textit{¿Cómo la educación afecta al crecimiento y a la desigualdad?}
\end{lnum}

Por tanto, la presencia de capital humano conduce a crecimiento a largo plazo en ausencia de progreso técnico exógeno. Con todo y siendo rigurosos, debemos tener en cuenta que la acumulación de capital humano difiere de la creación de conocimiento en la forma de progreso técnico. Si pensamos en el capital humano como las habilidades incorporadas en un trabajador, entonces el uso de estas habilidades en una actividad evita usarlas en otra actividad alternativa, lo que convierte al capital humano en un bien rival. Ya que la gente tiene derechos de propiedad sobre sus habilidades, el capital humano también es un bien excluible. En contraste, las ideas o el conocimiento puede ser no rival y, en algunas circunstancias también no excluible. 


\clearpage
\section{Segunda generación}
A partir de los años 90's, y habida cuenta de las limitaciones que suponían los modelos previos de crecimiento endogeneizado, la Teoría del Crecimiento Económico se centra en explicar las cómo la inversión en Investigación y Desarrollo, parte integral del comportamiento de las empresas, afecta y condiciona la senda del crecimiento económico. 

Como sabemos, la tecnología se define como el conjunto de conocimientos y procedimientos disponibles para transformar factores productivos en \textit{outputs}. Ahora bien, los procesos de I+D pueden también considerarse como un factor productivo propio que además genera externalidades, efecto desbordamiento, al conjunto de la economía. Por lo tanto, puede considerarse como un bien público, con características de no rivalidad y cierto grao de excluibilidad (atendiendo al sistema de patentes). No obstante, invertir en I+D también implica asumir unos costes, de hecho, unos costes fijos muy elevados, por lo que la empresa podría no invertir. Es necesario, para que esto no ocurra, establecer un sistema de protección intelectual que será modelizado a través del escenario de la competencia imperfecta (competencia monopolística).

Para incorporar la innovación tecnológica dentro de los programas de optimización, la literatura económica seguirá dos vías:
\begin{lnum}
    \item \textit{Aumento de las variedades}. La innovación tecnológica puede modelizarse a través del aumento de variedades en una economía que serán utilizados para producir el bien final;
    \item \textit{Aumento de la calidad}. La innovación tecnológica puede modelizarse a través del aumento de la calidad de los bienes de capital que serán utilizados en los procesos productivos (asemejable al proceso de \textit{destrucción creativa} de SHUMPETER).
\end{lnum}

\subsection{Crecimiento por acumulación de variedades de inputs. ROMER}
En este Tema desarrollaremos la primera aproximación a través de una versión simplificada del Modelo de ROMER (1990) con tres agentes.

\subsubsection{Productores finales}
Los productores finales utilizan como factores productivos el trabajo y los bienes intermedios (productos de innovación tecnológica protegidos con un sistema de patentes). 

\eqblock{
\begin{aligned}
    \text{Supuestos,}\;
    \begin{cases}
        \text{(1) Rendimietos Decrecientes respecto a variedades (\textit{input})}\\[0.5em]
        \text{(2) Rendimietos Constantes respecto a bienes finales}\\[0.5em]
        \text{(3) Alquilan bienes intermedios para producir:}\;Y_t=\underbrace{AL_t^{1-\alpha}\sum_{j=1}^Nx_{j,t}^\alpha}_{\text{Separable}}\\
        \hspace{2em}\text{Si }\forall j, x_{j,t}\text{ es la misma cantidad}\longrightarrow \overbrace{Y_t=AL_t^{1-\alpha}Nx^\alpha}^{\text{En equilibrio}}\\[0.5em]
        \text{(4) Mcdo. $L$ en competencia perfecta}\\[0.5em]
        \text{(5) Sistema de precios:}\;
        \begin{cases}
            \text{Bien intermedio:}\quad &P_{j,t}\\
            \text{Bien final:}\quad &P=1\\
        \end{cases}
    \end{cases}\\[1em]
    \text{Problema de maximización:}\quad
    \boxed{
    \begin{aligned}
        \max \sum_{t=0}^\infty \left[Y_t-w_tL_t-\sum_{j=1}^NP_{j,t}x_{j,t}\right]\\[0.5em]
    \end{aligned}
    }\\[0.5em]
    \text{CPO's}\;\;
    \begin{cases}
        w_t=PMg_L\\
        P_{j,t}=PMg_{x_{j,t}}\underbrace{\Longrightarrow}_{\text{Despejando}}\;
        \boxed{x_{j,t}^d=f(\overbrace{A,\alpha}^{+},\underbrace{P_{j,t}}_{-},L)}\quad \text{con }\underbrace{\varepsilon^d}_{\text{cte.}}=-\frac{1}{1-\alpha}
    \end{cases}
\end{aligned}
}{Problema de optimización: Productores finales. Modelo de ROMER (1990)}

\subsubsection{Empresas de bienes intermedios: I+D}
Las empresas de bienes intermedios operan en este modelo en eun mercado de competencia imperfecta. Son éstas las \textit{inventoras} de bienes de capital, para la cual necesitan esfuerzo y recursos en investigación y desarrollo. 

Además, para garantizar su rentabilidad, como ya hemos dicho, existe un marco legal que garantiza la protección de los Derechos de Propiedad Intelectual que convierte al productor en monopolista de su propia variedad. La producción de este bien intermedio dependerá, por tanto, de su $CMg$ y del $P_{j,t}$.

El problema del productor es en este caso maximizar el valor presente de los beneficios futuros de invertir en I+D, asumiendo que el trabajo es considerado constante y que el $CMg$ puede normalizarse a 1, tenemos:

\eqblock{
\begin{aligned}
    &\text{(1) Fijamos el precio:}\quad P_{j,t}=\underbrace{\frac{1}{\alpha}}_{\varepsilon^D}>1=CMg\\[0.5em]
    &\text{(2) Introducimos el precio:}\quad P_{j,t}^*\to x_{j,t}^d=f(A,\alpha,P_{j,t},L)\Longrightarrow x_{j,t}^{d^*}\\[0.5em]
    &\text{(3) Beneifcios descontados:}\quad \boxed{\int_0^{\infty}\pi\mathrm{d}t=\frac{\pi}{r}}\\[0.5em]
    &\text{(4) (Primera decisión) Decisión de inversión:}\\[1em]
    &\hspace{6em}\boxed{\int_0^{\infty}\pi\mathrm{d}t=\frac{\pi}{r}>0}\\[0.5em]
    &(\text{Se obtiene también la condición de entrada:}\quad \Phi=\frac{\pi}{r})
\end{aligned}
}{Problema del productor intermedio. Modelo de ROMER (1990)}

\subsubsection{Consumidores finales}
Por último, los consumidores finales se enfrentan a la decisión de invertir o consumir en cada periodo temporal, debiendo maximizar su utilidad fundamentada en el consumo total descontado o, equivalentemente, la senda de consumo intertemporal.

\eqblock{
\begin{aligned}
    &\boxed{
    \begin{aligned}
        \max_{c_t}U(0)=\int_{0}^\infty e^{\rho t}u_t(c_t)\mathrm{d}t\\[0.5em]
        \text{s.a.}\quad 
        \begin{cases}
            c_t+a_{t+1}\leq y_{t}\quad \text{(RPI)}\\[0.5em]
            \lim_{t\to \infty}a_t=0\\[0.5em]
            a_0,\quad \text{dado}
        \end{cases}
    \end{aligned}
    }\\
    &\text{CPO's}\Longrightarrow\text{Obtenemos:}\;\boxed{\gamma_c=\frac{1}{\theta}(r-\rho)\sim\frac{1}{\theta}(\frac{\pi}{\Phi}-\rho)}\Longrightarrow\boxed{\gamma_c=f\left(\overbrace{\frac{1-\alpha}{\alpha}}^{+}, \underbrace{\Phi}_{-}\right)}
\end{aligned}
}{Problema de los hogares. Modelo de ROMER (1990)}

Por tanto, obtenemos que la senda de crecimiento del consumo es constante en el equilibrio y mantiene una relación positiva con el markup ($\mu$) de las empresas de innovación en la fijación de precios, e inversa con las barreras de entrada existentes en dicho sector. 

Este último fenómeno se debe a los Efectos Escala, ya que cuando hay mayor población, la senda de crecimiento del consumo es mayor ya que, como la tecnología es un bien no rival y la proporción de recursos de capital dedicados al sector I+D es constante, cuanto más aumente la población menores serán los costes del sector I+D en términos per cápita.

\subsubsection{Valoración en términos de Política Económica}
Por tanto, como vemos, ROMER logra endogeneizar el progreso técnico como parte de su modelo introduciendo competencia imperfecta, y es la invención de nuevas variedades de bienes intermedios el motor que nos permite crecer de forma sostenida, a una tasa constante e igual para todas la variables.

\paragraph{Planificador benevolente}
Además, una de las implicaciones más importantes en términos de Política Económica es que la solución del planificador y la solución descentralizada no coinciden, puesto que nos encontramos en un supuesto de competencia imperfecta (fallos de mercado) y externalidades. 

El precio fijado por las empresas innovadoras es superior al coste marginal, lo que reduce la cantidad demandada por la empresa del bien final (y por ende su producción). Esto tiene dos efectos:
\begin{la}
    \item \textbf{Efecto nivel}. Al incorporar menos capital se reduce la producción del output; y,
    \item \textbf{Efecto tasa de crecimiento}. Al vender una menor cantidad se reduce la entrada de ideas que generen nuevas variedades. 
\end{la}

Ahora bien, esto no supone que haya que eliminar las patentes, pues sin éstas directamente no habría producción de innovaciones. ¿Qué soluciones pueden existir?
\begin{lnum}
    \item Dar subsidios a la investigación: de forma que se reduzcan los costes fijos a los que las empresas innovadoras tienen que hacer frente, se incentivaría que más empresas entrasen a innovar y que el crecimiento fuese entonces superior. El problema de esta medida es que no consigue corregir la distorsión, pues seguirían comprando un producto con un precio superior al coste marginal. Esta medida corregiría la tasa de crecimiento del PIB pero no el nivel; o,
    \item Subsidiar las compras de productos intermedios: consistiría en proporcionar subsidios (de suma fija) a las empresas de bien final para que adquieran la cantidad competitiva de bienes intermedios, es decir, que aunque los productores perciban un precio superior al coste marginal, las empresas que adquieren esos bienes intermedios solo paguen un precio igual al coste marginal. Así, se eliminaría toda distorsión (subsidios de suma fija) y se permitiría alcanzar el óptimo resolviendo tanto el problema en niveles como el problema en tasa. 
\end{lnum}
+ 
 
2º) Los modelos de bienes intermedios antiguos siguen siendo válidos: de la forma en la que hemos agregado los bienes intermedios hemos visto que todos son simétricos, es decir, todos tienen un mismo impacto en el output y por tanto son igual de deseados por la empresa. Esto lleva, no obstante, a que la maquinaria inventada hace 10 años sea igual de buena que la actual, lo que es poco realista.
Por este motivo surgió en la literatura otra corriente de modelos de 2ª generación, con la misma esencia que el modelo de Romer aquí presentado, pero que consideran que la innovación no es la creación de nuevas variedades, sino la creación de mejores calidades, como el modelo de Aghion y Howit. Esto es importante, pues supone que la creación de una variedad nueva expulsa del mercado a la vieja, que ya no es útil. Estos modelos son especialmente relevantes al considerar a la innovación como un proceso de destrucción, aplicando esa idea de la “destrucción creativa” de Schumpeter.= Esto tiene enormes implicaciones en términos de política económica, pues estos modelos se caracterizan precisamente por tener una tasa de innovación excesivamente elevada.
La razón es que los inventores no internalizan las pérdidas que están generando de la variedad expulsada del mercado. El planificador sí que lo tiene en cuenta. Esto lleva a un exceso de inversión y a un crecimiento excesivo, lo que abre la puerta a consideraciones como la necesidad de compensar a los perdedores del proceso innovador. 


\eqblock{
\begin{aligned}
    
\end{aligned}
}




\clearpage
\section{Conclusión} 




\subsection{Recapitulación}



\subsection{Extensiones y relación con otras partes del Temario}



\subsection{Opinión}

\subsubsection{Idea final (Salida o cierra)}

\clearpage
\pagenumbering{gobble}
\MyBibliography


\MyBibItem{DAWID y GATTI}{Chapter 2 - Agent-Based Macroeconomics}{2018}{https://doi.org/10.1016/bs.hescom.2018.02.006}


% --- Anexos ---
\clearpage
\anexo{\textbf{Preguntas Test}}
%\input{\TemaCode/questions.tex} 



\clearpage
\anexo{Historia del desarrollo teórico}\label{anx:historia}


\end{document}